\section*{The SPADI Questionnaire: A Worked Example}

This appendix applies the full item screening procedure of Section~\ref{sec:it_sec} to a real dataset: the Shoulder Pain and Disability Index (SPADI).

\subsection{Step 0: Getting Started}
Before screening, we inspect the data and perform any necessary preprocessing. We refer to this preparation as Step 0, although it is not a formal step of the screening procedure.


The Shoulder Pain and Disability Index comprises $13$ items: five assess pain and eight assess disability. Responses range from $0$ to $10$, with higher scores indicating greater pain or disability. The dataset also contains two covariates, \texttt{gender} and \texttt{over60}.

Of the $229$ respondents, $213$ have complete item responses. Neither covariate has missing values among these $213$ respondents, so all are included in the analysis. Following \citet{christensen-et-al-2019}, we analyze the pain and disability subscales separately and recode the response categories from $(0,1,2,3,4,5,6,7,8,9,10)$ to $(0,0,1,1,2,2,3,3,4,4,5)$. The resulting datasets have dimensions $213\times 7$ for pain and $213\times 10$ for disability, including the two covariates in each.


\subsection{Step 1}
Both subscales pass the M1 and M2 checks (results not shown). The covariate \texttt{over60} fails M3 for both subscales. This would call further screening into question if age were treated as a criterion variable. We do not make that assumption here, so we proceed.


\subsection{Step 2}
We control the false discovery rate using the Benjamini--Hochberg procedure. The initial screen finds no DIF in the disability subscale. In the pain subscale, item $P1$ shows DIF with respect to \texttt{over60} ($\gamma=-0.43$, $p<0.001$). Both subscales show evidence of LD. For disability, the pair $(D4,D5)$ is significant in both conditioning directions. For pain, $(P3,P4)$ is significant in both directions, whereas $(P1,P4)$ is significant only when conditioning on $R_{P4}$.

\subsection{Step 3}
\subsubsection{Step 3(a)}
After the inclusion algorithm, we find the following pairs of genuine LD in the two subscales:
\begin{table}[H]
    \centering
    \begin{minipage}{0.45\textwidth}
        \centering
        \begin{tabular}{c c c}
            & Disability subscale & \\
            \hline
            Item1 & Item2 & partial $\gamma$ \\ 
            \hline 
            $D4$ & $D5$ & $0.425$ \\
            \hline 
        \end{tabular}
        \caption{Genuine LD pairs.}
        \label{tab:disability_partial_gamma}
    \end{minipage}
    \hfill
    \begin{minipage}{0.45\textwidth}
        \centering
        \begin{tabular}{c c c}
            & Pain subscale & \\
            \hline
            Item1 & Item2 & partial $\gamma$ \\ 
            \hline 
            $P3$ & $P4$ & $0.467$ \\
            \hline 
        \end{tabular}
        \caption{Genuine LD pairs.}
        \label{tab:pain_partial_gamma}
    \end{minipage}
\end{table}

\subsubsection{Steps 3(b) and 3(c)}
Because no DIF was detected in the disability subscale, Steps 3(b) and 3(c) apply only to the pain subscale. We use asymptotic $p$-values, consistent with the rest of the analysis. The initial screen identifies only one DIF relationship in this subscale, so the elimination procedures return that relationship and terminate.
\begin{table}[H]
    \centering
    \begin{tabular}{c c c c c c c c c}
         item & source & partial $\gamma_1$ & $\mathcal{C}_1$ & $p_1$ & partial $\gamma_2$ & $\mathcal{C}_2$ & $p_2$ & conclusion \\
         \hline
          $P1$   &  over60 &     $-0.430$      &      Score & $0.0009$ &  $-0.430$   &  Score  &  $0.0009$& DIF\\
          \hline
    \end{tabular}
    \caption{Test for genuine DIF. Here, $\mathcal{C}_i$ denotes the conditioning set for the corresponding partial gamma coefficient.}
    \label{tab:step_3bc_ex}
\end{table}

\subsection{Step 4}
For the disability subscale, both covariates remain significantly associated with the score after conditioning on the other and are classified as supporting criterion validity. Neither covariate meets this criterion for the pain subscale.
\begin{table}[H]
    \centering
    \begin{tabular}{c|c c c c c}
    \hline 
         Subscale  & covariate & partial $\gamma$ & $p$ & $\mathcal{C}$ & S.C.V. \\
         \hline
         Disability &  gender &  $0.248$ &  $0.0007$               &   over60 &    TRUE \\
        & over60 & $0.190$  & $0.027$ & gender & TRUE \\
        \hline
        Pain & gender & $0.080$ & $0.115$  & over60 &   FALSE \\
        &  over60 & $0.029$ & $ 0.860$ & gender&  FALSE \\
        \hline
    \end{tabular}
    \caption{$\mathcal{C}$ is the conditioning set; S.C.V.\ abbreviates ``supports criterion validity''.}
    \label{tab:placeholder}
\end{table}


\subsection{Step 5}
We construct the graphs for the two subscales from the retained screening results.
\begin{figure}[ht]
\centering

% -------------------- (a) IRT graph --------------------
\begin{minipage}[t]{0.48\textwidth}
\centering
\begin{tikzpicture}[scale=0.68,transform shape,
  every node/.style={font=\scriptsize}]

% Items, placed on a left arc
\node[nodecircle] (D1) at (-3.2,  2.8) {D1};
\node[nodecircle] (D2) at (-4.0,  2.0) {D2};
\node[nodecircle] (D3) at (-4.4,  1.0) {D3};
\node[nodecircle] (D4) at (-4.5,  0.0) {D4};
\node[nodecircle] (D5) at (-4.5, -1.0) {D5};
\node[nodecircle] (D6) at (-4.4, -2.0) {D6};
\node[nodecircle] (D7) at (-4.0, -3.0) {D7};
\node[nodecircle] (D8) at (-3.2, -3.8) {D8};

% Latent variable in the middle
\node[nodecircle] (theta) at (0, -0.5) {$\theta$};

% Exogenous variables on the right
\node[nodecircle] (gender) at (3.8, 1.3) {gender};
\node[nodecircle] (over60) at (3.8,-2.3) {over60};

% Measurement edges
\draw[->,thick] (theta) -- (D1);
\draw[->,thick] (theta) -- (D2);
\draw[->,thick] (theta) -- (D3);
\draw[->,thick] (theta) -- (D4);
\draw[->,thick] (theta) -- (D5);
\draw[->,thick] (theta) -- (D6);
\draw[->,thick] (theta) -- (D7);
\draw[->,thick] (theta) -- (D8);

% Score associations
\draw[dif] (gender) -- (theta);
\draw[dif] (over60) -- (theta);

% Local dependence
\draw[ld] (D4) -- (D5);

\end{tikzpicture}

\vspace{0.3em}
\textbf{(a)}
\end{minipage}
\hfill
% -------------------- (b) Moralized graph --------------------
\begin{minipage}[t]{0.48\textwidth}
\centering
\begin{tikzpicture}[scale=0.68,transform shape,
  every node/.style={font=\scriptsize}]

% Items, placed on a left arc / clique
\node[nodecircle] (D1) at (-3.2,  2.8) {D1};
\node[nodecircle] (D2) at (-4.0,  2.0) {D2};
\node[nodecircle] (D3) at (-4.4,  1.0) {D3};
\node[nodecircle] (D4) at (-4.5,  0.0) {D4};
\node[nodecircle] (D5) at (-4.5, -1.0) {D5};
\node[nodecircle] (D6) at (-4.4, -2.0) {D6};
\node[nodecircle] (D7) at (-4.0, -3.0) {D7};
\node[nodecircle] (D8) at (-3.2, -3.8) {D8};

% Score in the middle
\node[nodecircle] (S) at (0, -0.5) {$S$};

% Exogenous variables on the right
\node[nodecircle] (gender) at (3.8, 1.3) {gender};
\node[nodecircle] (over60) at (3.8,-2.3) {over60};

% Score-item edges
\draw[scoreedge] (D1) -- (S);
\draw[scoreedge] (D2) -- (S);
\draw[scoreedge] (D3) -- (S);
\draw[scoreedge] (D4) -- (S);
\draw[scoreedge] (D5) -- (S);
\draw[scoreedge] (D6) -- (S);
\draw[scoreedge] (D7) -- (S);
\draw[scoreedge] (D8) -- (S);

% Item-score moralization clique
\foreach \i/\j in {
D1/D2,D1/D3,D1/D4,D1/D5,D1/D6,D1/D7,D1/D8,
D2/D3,D2/D4,D2/D5,D2/D6,D2/D7,D2/D8,
D3/D4,D3/D5,D3/D6,D3/D7,D3/D8,
D4/D6,D4/D7,D4/D8,
D5/D6,D5/D7,D5/D8,
D6/D7,D6/D8,
D7/D8}
{
  \draw[thin,gray!55] (\i) -- (\j);
}

% Local dependence, emphasized on top of the moralization edge
\draw[ld] (D4) -- (D5);

% Score associations
\draw[scoreedge] (gender) -- (S);
\draw[scoreedge] (over60) -- (S);

% Moralized parent edge
\draw[ld] (gender) -- (over60);

\end{tikzpicture}

\vspace{0.3em}
\textbf{(b)}
\end{minipage}

\caption{
Disability subscale graphs: (a) IRT graph; (b) corresponding moralized graph.}\label{fig-disabilitygraph}
\end{figure}


\begin{figure}[ht]
\centering

% -------------------- (a) IRT graph --------------------
\begin{minipage}[t]{0.48\textwidth}
\centering
\begin{tikzpicture}[scale=0.78,transform shape,
  every node/.style={font=\scriptsize}]

% Items, placed on a left arc
\node[nodecircle] (P1) at (-3.4,  2.0) {P1};
\node[nodecircle] (P2) at (-4.1,  1.0) {P2};
\node[nodecircle] (P3) at (-4.2,  0.0) {P3};
\node[nodecircle] (P4) at (-4.1, -1.0) {P4};
\node[nodecircle] (P5) at (-3.4, -2.0) {P5};

% Latent variable in the middle
\node[nodecircle] (theta) at (0, 0) {$\theta$};

% Exogenous variable on the right
\node[nodecircle] (over60) at (3.7, 0.8) {over60};
\node[nodecircle] (age) at (3.7, -0.8) {gender};

% Measurement edges
\draw[->,thick] (theta) -- (P1);
\draw[->,thick] (theta) -- (P2);
\draw[->,thick] (theta) -- (P3);
\draw[->,thick] (theta) -- (P4);
\draw[->,thick] (theta) -- (P5);

% DIF edge
\draw[dif] (over60) -- (P1);

% Local dependence
\draw[ld] (P3) -- (P4);

\end{tikzpicture}

\vspace{0.3em}
\textbf{(a)}
\end{minipage}
\hfill
% -------------------- (b) Moralized graph --------------------
\begin{minipage}[t]{0.48\textwidth}
\centering
\begin{tikzpicture}[scale=0.78,transform shape,
  every node/.style={font=\scriptsize}]

% Items, placed on a left arc / clique
\node[nodecircle] (P1) at (-3.4,  2.0) {P1};
\node[nodecircle] (P2) at (-4.1,  1.0) {P2};
\node[nodecircle] (P3) at (-4.2,  0.0) {P3};
\node[nodecircle] (P4) at (-4.1, -1.0) {P4};
\node[nodecircle] (P5) at (-3.4, -2.0) {P5};

% Score node in the middle; named S in TikZ, displayed as #
\node[nodecircle] (Score) at (0, 0) {$S$};

% Exogenous variable on the right
\node[nodecircle] (over60) at (3.7, 0.8) {over60};
\node[nodecircle] (age) at (3.7, -0.8) {gender};

% Score-item edges
\draw[scoreedge] (Score) -- (P1);
\draw[scoreedge] (Score) -- (P2);
\draw[scoreedge] (Score) -- (P3);
\draw[scoreedge] (Score) -- (P4);
\draw[scoreedge] (Score) -- (P5);

% Item-score moralization clique
\foreach \i/\j in {
P1/P2,P1/P3,P1/P4,P1/P5,
P2/P3,P2/P4,P2/P5,
P3/P5,
P4/P5}
{
  \draw[thin,gray!55] (\i) -- (\j);
}

% Local dependence, emphasized on top of the moralization edge
\draw[ld] (P3) -- (P4);

% DIF edge in the moralized graph
\draw[ld] (over60) -- (P1);

% Moralized parent edge
\draw[ld] (Score) -- (over60);

\end{tikzpicture}

\vspace{0.3em}
\textbf{(b)}
\end{minipage}

\caption{
Pain subscale graphs: (a) IRT graph; (b) corresponding moralized graph.
}\label{fig-paingraph}
\end{figure}


\subsection{Step 6}
Step 6 retests the retained DIF and LD edges using the conditioning sets implied by the graphs. For these two subscales, the LD and DIF structures make the additional tests redundant. The final graphs are therefore those reported in Step 5 (Figures~\ref{fig-disabilitygraph} and~\ref{fig-paingraph}).


